\section{Outlook: Discovery as Search, Language as Realization}
\label{sec:outlook}

A scope statement first. Nothing in this paper shows that large
language models cannot originate scientific questions \emph{in
principle}; a single prompting strategy against a single model family
in a single domain cannot support that claim, and we do not make it.
What the data do support is narrower and more useful: \emph{bare
prompting does not reliably perform problem discovery}, and the
components of the systems that do perform better can be named.

\paragraph{Three capabilities the experiments separate.} Read together,
the decomposition (Section~\ref{sec:foresight:abc}), the anchoring
rubric (Section~\ref{sec:foresight:saf}), and the temporal stress test
(Section~\ref{sec:foresight:stress}) distinguish three things that
``asking good questions'' conflates. First, a \emph{topic prior}:
knowing what a field is likely to work on next. This is what LLM-only
generation exhibits---near-ceiling engagement in every era (92--98\%),
questions that name object classes rather than objects (5\% anchoring
against 95--100\% for the structural pipelines), and
specificity-adjusted foresight within a few points of the
random-template floor. A topic prior is genuinely predictive of where
attention flows, and genuinely cheap: it requires no memory of the
future, which is why it survives the training boundary unchanged.
Second, \emph{structural problem discovery}: locating specific,
contestable configurations of existing evidence---a claim, a
counter-claim, a method dependency, an untested premise. The
weight-free pipeline C is the clean witness that this capability does
not reside in model weights: with no language model anywhere, it finds
questions the future substantively engages at every cutoff, and it
beats LLM-only prompting on resolution. Third,
\emph{articulation}: turning a detected configuration into a question a
scientist would recognize as askable. This is where the LLM earns its
place---over identical evidence structures, LLM verbalization roughly
doubles resolution and quintuples refutations relative to a fixed
template---and it is a capability the stress test shows to be
era-robust rather than memorized.

\paragraph{The architecture this implies.} These results point away
from ``make the model smarter and ask it for a hundred ideas'' and
toward a division of labour:
\begin{quote}
\emph{machine-scale structural search over the evidence space
$\;\rightarrow\;$ candidate scientific tensions $\;\rightarrow\;$ LLM
articulation $\;\rightarrow\;$ testable questions.}
\end{quote}
The asymmetry that motivates it is quantitative. A literature of
thousands of papers yields tens of thousands of claims and a
combinatorially larger space of claim pairs, evidence paths, and
method dependencies---far beyond what any single reader, human or
prompted model, holds in attention at once, but squarely within what a
machine can sweep in parallel. Under this framing the LLM is never
asked to conjure novelty from nothing; it is asked to do what it
demonstrably does well---local semantic understanding, claim
extraction, relation judgment, and finally phrasing---while the search
system carries the burden of combinatorial exploration. Scientific
question discovery becomes a \emph{computable search problem} over an
explicit representation of what the literature claims, and the
interesting engineering question shifts from prompting to
representation and search: what to index, which configurations to
enumerate, and how to rank them.

\paragraph{The next measurable question.} Ranking is where this
benchmark and that architecture meet. Every question in the released
instances carries its generating signal---tension type, object,
method-dependency, source claims---and its measured fate. That pairing
makes a new question answerable: \emph{which structural patterns most
often lead to questions the future answers, advances, or
refutes?} Our own probe is deliberately crude---object co-mention plus
stance cues, visibly below the human-gated graph in precision---so the
headroom is real: contradiction typing, dataset-dependency detection,
and archival-data availability are all candidate features for a
learned prior over the tension space. We flag the honest status of all
of this: an interpretation consistent with our data, not an
established causal account---and the frozen prospective instance
(Section~\ref{sec:conclusion}) is the experiment that will test it
without any of this paper's retrospective caveats. Learning that prior
from backtested outcomes, and validating it prospectively, is the
natural next paper.
